\section{Benchmarking Details}
\label{app:benchmarking}

This appendix records the benchmark methodology and supporting code behind
Section~\ref{sec:evaluation}.  The intent is reproducibility rather than a new
claim beyond the main paper: the artifact reports both successful derivations
and unsupported cases, and it compares generated serializers against handwritten
Rocq baselines extracted through the same pipeline.

\subsection{Applicability Benchmark}

The artifact scans installed Corelib and Stdlib sources for inductive-like
declarations, probes each candidate in isolation, records diagnostics, and then
compiles a single benchmark file containing all successful derivations.  This
workflow is deliberately conservative: unsupported declarations remain visible
as categorized failures.  Declarations that are visible in source but are not
closed derivation targets, such as declarations inside functor bodies or module
signatures, are recorded separately rather than discarded.

\begin{table}
  \centering
  \caption{Corelib/Stdlib applicability and failure categories.}
  \label{tab:app-detailed-applicability}
  \small
  \begin{minipage}[t]{0.47\linewidth}
    \centering
    \begin{tabular}{p{0.72\linewidth}r}
      \toprule
      Metric                       & Count                        \\
      \midrule
      Discovered declarations      & \JsonifiableDiscovered       \\
      Recorded, not probed         & \JsonifiableSkipped          \\
      Probed declarations          & \JsonifiableProbed           \\
      Successful probes            & \JsonifiableProbeSuccesses   \\
      Probe failures               & \JsonifiableProbeFailures    \\
      Probe timeouts               & \JsonifiableProbeTimeouts    \\
      Timed successful derivations & \JsonifiableTimedDerivations \\
      \bottomrule
    \end{tabular}
  \end{minipage}\hfill
  \begin{minipage}[t]{0.47\linewidth}
    \centering
    \begin{tabular}{p{0.78\linewidth}r}
      \toprule
      Failure category & Count \\
      \midrule
      \JsonifiableFailureRows
      \bottomrule
    \end{tabular}
  \end{minipage}
\end{table}

The category names in Table~\ref{tab:app-detailed-applicability} are a triage
device, not a claim that every failed probe is a tool defect.  Some declarations
are clear non-targets for serialization: propositions are erased by extraction,
functions have no canonical JSON representation, sort-valued fields contain
types rather than values, and coinductive values may be infinite.  The genuine
current limitations are concentrated around dependent data, where decoding must
reconstruct values whose types depend on earlier fields or external indices.
Table~\ref{tab:failure-examples} gives one representative declaration for each
category and the rationale used by the artifact.

\begin{table}
  \centering
  \caption{Representative failure examples and classification rationale.}
  \label{tab:failure-examples}
  \scriptsize
  \setlength{\tabcolsep}{3pt}
  \begin{tabular}{p{0.43\linewidth}p{0.53\linewidth}}
    \toprule
    Category and example & Rationale                                                                                                  \\
    \midrule
    \texttt{prop-target-not-supported}: \texttt{or}
                         & The target lives in \texttt{Prop}; propositions are erased at extraction and are not runtime data.         \\
    \texttt{dependent-field-not-supported}: \texttt{sig}
                         & A later field depends on an earlier witness; decoding must reconstruct a value and its dependent type.     \\
    \texttt{indexed-family-not-supported}: \texttt{reflect}
                         & The declaration is indexed by values; decoding requires recovering the index or fixing it in the instance. \\
    \texttt{function-field-not-supported}: \texttt{implies}
                         & Arbitrary functions have no canonical data representation in JSON or ordinary serializers.                 \\
    \texttt{proof-field-not-supported}: \texttt{sumbool}
                         & A computational wrapper carries proof evidence; decoding requires proof erasure or reconstruction.         \\
    \texttt{sort-field-not-supported}: \texttt{Tlist}
                         & Constructor fields contain sorts such as \texttt{Type}; these are types, not JSON values.                  \\
    \texttt{coinductive-not-supported}: \texttt{Stream}
                         & Potentially infinite data needs an explicit depth-bounded or lazy serialization contract.                  \\
    \bottomrule
  \end{tabular}
\end{table}

The non-target group is not a design choice particular to
\texttt{rocq-json}: there is no runtime proposition to serialize after
extraction; arbitrary functions do not have a canonical finite data
representation; sorts such as \texttt{Type} and \texttt{Set} are not ordinary
values; and coinductive data needs a separate bounded or lazy contract.  These
account for 158 probe failures.  The remaining large categories are dependent
fields and indexed families, 69 cases in total.  Those are real future work:
the JSON decoder would have to recover enough information to reconstruct a
dependent value, or the deriver would need a fixed-index or sigma-style
encoding.  The three proof-field cases sit on the same boundary: the outer type
is computational, but successful decoding would require erasing or rebuilding
proof evidence that extraction normally removes.

The \JsonifiableSkipped{} recorded-but-not-probed declarations are reported in
the artifact as well: \JsonifiableSkippedParameterizedModule{} occur inside
parameterized modules or functor bodies, such as the raw tree declaration in
\texttt{FMapAVL}, and
\JsonifiableSkippedModuleType{} occur inside module types or signatures, such
as the tree declaration in \texttt{MSetGenTree.Ops}.  These are genuine source
declarations, but not closed global references to which the deriver can be
applied directly \emph{outside the functor}.

Excluding the \JsonifiableSkipped{} functor/signature declarations that are not
valid deriver inputs and the 158 non-target declarations, the tool successfully
derives instances for \JsonifiableProbeSuccesses{} of 153 semantically eligible
declarations, roughly 53\% coverage on the well-defined target fragment, with
the remaining gap concentrated in dependent and indexed types.

\subsection{Derivation Time}

The combined successful benchmark records Rocq's \texttt{Time} output for each
derived instance.  In the reported run, the combined benchmark compiled
successfully with wall-clock time \JsonifiableBenchmarkWallSeconds{} seconds;
the sum of Rocq-reported derivation transactions was
\JsonifiableDerivationSumSeconds{} seconds.  The mean derivation time was
\JsonifiableDerivationMeanSeconds{} seconds, the median was
\JsonifiableDerivationMedianSeconds{} seconds, the 90th percentile was
\JsonifiableDerivationPNineZeroSeconds{} seconds, and the maximum was
\JsonifiableDerivationMaxSeconds{} seconds.

\begin{table}
  \centering
  \caption{Slowest successful derivations in the broad benchmark.}
  \label{tab:slowest-derivations}
  \begin{tabular}{l @{\hspace{0.25cm}} lr}
    \toprule
    Declaration & Kind & Seconds \\
    \midrule
    \JsonifiableSlowestRows
    \bottomrule
  \end{tabular}
\end{table}

The slowest success is \texttt{byte}, a 256-constructor Corelib enumeration.
The remaining cost is dominated by Rocq-side derivation and proof construction.
After extraction, the corresponding dispatch path is close to the handwritten
Rocq baseline used in the extracted-code benchmark.

\subsection{Extracted Code Performance}

The extraction benchmark compares generated definitions against handwritten
Rocq definitions after both have been extracted to OCaml.  This keeps the
comparison in the same source language and extraction pipeline.  The benchmark
covers several shapes: a 256-constructor enumeration, a tuple-like constructor,
a flat record, two mixed sums with nullary and field-bearing constructors, and a
recursive binary tree.  Each benchmark was run
\JsonifiableExtractionGeneratedRuns{} times.  The enum benchmark serializes and
deserializes every constructor; the other benchmarks run many iterations over
representative values.  The handwritten baselines were written by the authors in
Rocq and extracted through the same pipeline; they were not independently
optimized.

\subsection{Round-Trip Benchmark Rocq Code}

The extracted-code benchmarks in Table~\ref{tab:extraction} are defined in the
extraction benchmark source.  The generated cases use
\texttt{Elpi derive.jsonifiable}; the primed instances below are the handwritten
Rocq baselines extracted and timed by the same harness.

\noindent\textbf{Benchmark Types}

\lstinputlisting[
  language=Coq,
  firstline=26,
  lastline=45,
  breaklines=true,
  basicstyle=\ttfamily\scriptsize
]{../extracting/JSON_Derive_Tests.v}

\lstinputlisting[
  language=Coq,
  firstline=68,
  lastline=76,
  breaklines=true,
  aboveskip=0pt,
  belowskip=0pt,
  basicstyle=\ttfamily\scriptsize
]{../extracting/JSON_Derive_Tests.v}

\lstinputlisting[
  language=Coq,
  firstline=165,
  lastline=168,
  breaklines=true,
  aboveskip=0pt,
  belowskip=0pt,
  basicstyle=\ttfamily\scriptsize
]{../extracting/JSON_Derive_Tests.v}

\lstinputlisting[
  language=Coq,
  firstline=425,
  lastline=458,
  breaklines=true,
  aboveskip=0pt,
  belowskip=0pt,
  basicstyle=\ttfamily\scriptsize
]{../extracting/JSON_Derive_Tests.v}

\noindent\textbf{Handwritten Rocq Baselines}

\lstinputlisting[
  language=Coq,
  firstline=257,
  lastline=383,
  breaklines=true,
  aboveskip=0pt,
  belowskip=0pt,
  basicstyle=\ttfamily\scriptsize
]{../extracting/JSON_Derive_Tests.v}

\lstinputlisting[
  language=Coq,
  firstline=461,
  lastline=471,
  breaklines=true,
  aboveskip=0pt,
  belowskip=0pt,
  basicstyle=\ttfamily\scriptsize
]{../extracting/JSON_Derive_Tests.v}
$$\vdots$$
\lstinputlisting[
  language=Coq,
  firstline=590,
  lastline=601,
  breaklines=true,
  aboveskip=0pt,
  belowskip=0pt,
  basicstyle=\ttfamily\scriptsize
]{../extracting/JSON_Derive_Tests.v}
$$\vdots$$
\lstinputlisting[
  language=Coq,
  firstline=720,
  lastline=724,
  breaklines=true,
  aboveskip=0pt,
  belowskip=0pt,
  basicstyle=\ttfamily\scriptsize
]{../extracting/JSON_Derive_Tests.v}

\noindent\textbf{Extraction Command}

\lstinputlisting[
  language=Coq,
  firstline=726,
  lastline=734,
  breaklines=true,
  aboveskip=0pt,
  belowskip=0pt,
  basicstyle=\ttfamily\scriptsize
]{../extracting/JSON_Derive_Tests.v}
